\section{Composition and Cross-Diagram Links}
\label{sec:composition}

Triton's distinguishing capability is putting whole diagrams together on one
canvas and drawing connections \emph{between} them. Two features cooperate: the
\texttt{poster} grid and cross-diagram links.

\subsection{Posters}

A \texttt{poster} arranges child diagrams in a grid. A cell may hold \emph{any}
registered diagram kind --- the cell kind is dispatched through the same registry
the top-level pipeline uses (\texttt{src/diagrams/poster/}), so a poster can embed
a plan tree next to a B-tree next to a slotted page. Cells flow left-to-right and
wrap at the column count; a cell may span columns or rows:

\begin{lstlisting}[language=triton]
poster "Cell Spanning"
    columns 3

    cell hero "spans all 3 columns" [3] :: plan
        plan
            Hash Join {rows: 980}
                Seq Scan orders {rows: 10000}
    end

    cell idx "spans 2" [2] :: btree
        btree order 3 insert 10 20 30 40 50 60
    end
    cell k "fits the 3rd" :: stat
        13 | kinds
    end
\end{lstlisting}

\ourfig{spanning}{A poster demonstrating cell spanning: a plan tree spanning all
three columns, a B-tree spanning two, and single-width cells flowing into the
remaining slots.}

\texttt{[n]} spans \texttt{n} columns; \texttt{[c x r]} spans \texttt{c} columns
by \texttt{r} rows. Placement is occupancy-aware: a row-spanning cell reserves the
rows it covers, so later cells flow past it without overlap.

\subsection{Cross-diagram links}

Because each diagram's \texttt{layout} returns a node-anchor registry
(\S\ref{sec:diagram-contract}), the poster layer can address a node \emph{inside}
a child diagram as \texttt{cellId.nodeId} and route an arrow to a node in a
different cell:

\begin{lstlisting}[language=triton]
    link q.n3  --> idx.n0 "scan uses index"
    link idx.n0 --> pg.slot0 "leaf points to page"
\end{lstlisting}

The poster merges child anchor registries with their cell id as a path prefix,
then routes the link with the same slot-aware connector used everywhere else. This
is what lets, for example, an ``SQL engine anatomy'' poster show a query plan, a
B-tree index, and a storage page as one connected story.
